\documentclass[12pt, a4paper]{article}
\usepackage[utf8]{inputenc}
\usepackage[T1]{fontenc}
\usepackage{amsmath, amssymb}
\usepackage{geometry}
\usepackage{booktabs}
\usepackage{titlesec}
\usepackage{xcolor}
\usepackage{mdframed}
\usepackage{enumitem}
\usepackage{array}
\usepackage{parskip}

\geometry{margin=2.5cm, top=3cm, bottom=3cm}

\definecolor{ballblue}{RGB}{0, 82, 147}
\definecolor{lightgray}{RGB}{245, 245, 245}
\definecolor{darkgray}{RGB}{80, 80, 80}
\definecolor{accentgray}{RGB}{180, 180, 180}

\titleformat{\section}
  {\large\bfseries\color{ballblue}}{}
  {0em}{}[\color{ballblue}\titlerule]
\titlespacing{\section}{0pt}{18pt}{8pt}

\mdfdefinestyle{destaque}{
  backgroundcolor=lightgray,
  linecolor=ballblue, linewidth=1.2pt,
  leftline=true, rightline=false, topline=false, bottomline=false,
  innerleftmargin=10pt, innerrightmargin=10pt,
  innertopmargin=6pt, innerbottommargin=6pt,
  skipabove=4pt, skipbelow=4pt,
}

% =============================================================================
\begin{document}

% ── Capa ─────────────────────────────────────────────────────────────────────
\begin{center}
  {\color{ballblue}\rule{\linewidth}{2pt}}\\[8pt]
  {\LARGE\bfseries Modelo de Programa\c{c}\~{a}o Linear}\\[4pt]
  {\large\color{darkgray} Planejamento Macrosc\'{o}pico --- Ball Corporation}\\[8pt]
  {\color{ballblue}\rule{\linewidth}{2pt}}
\end{center}

\vspace{10pt}
\begin{center}
  \small\color{darkgray}
  \textbf{1.~Modelo de minimiza\c{c}\~{a}o de custos}
  \enspace\textbullet\enspace
  \textbf{2.~Conjuntos}
  \enspace\textbullet\enspace
  \textbf{3.~Par\^{a}metros}
  \enspace\textbullet\enspace
  \textbf{4.~Fun\c{c}\~{a}o Auxiliar}
  \enspace\textbullet\enspace
  \textbf{5.~Vari\'{a}veis de Decis\~{a}o}
  \enspace\textbullet\enspace
  \textbf{6.~Fun\c{c}\~{a}o Objetivo}
\end{center}
\vspace{4pt}
{\color{accentgray}\rule{\linewidth}{0.5pt}}

% ============================================================
\section{Modelo de minimiza\c{c}\~{a}o de custos}

O modelo busca \textbf{minimizar o custo total de opera\c{c}\~{a}o} da malha industrial da Ball
Corporation, contemplando tr\^{e}s componentes de custo: produ\c{c}\~{a}o, estocagem e transfer\^{e}ncia
log\'{i}stica entre plantas. O modelo \'{e} estritamente cont\'{i}nuo (sem vari\'{a}veis bin\'{a}rias ou
inteiras) e cobre um horizonte de 15 dias.

\[
  \boxed{
    \begin{aligned}
      \min \; Z \;=\;\;
      & \underbrace{\sum_{p \in P}\sum_{l \in L_p}\sum_{i \in I}\sum_{t \in T}
          k_{i,p,l} \cdot x_{i,p,l,t}}_{\text{custo de produ\c{c}\~{a}o}}
        \;+\;
        \underbrace{\sum_{p \in P}\sum_{l \in L_p}\sum_{i \in I}\sum_{t \in T}
          v_{i,p,l} \cdot e_{i,p,l,t}}_{\text{custo de estoque}} \\[14pt]
      & +\;
        \underbrace{\sum_{i \in I}\sum_{o \in P}\sum_{\substack{d \in P \\ d \neq o}}\sum_{t \in T}
          ct_{o,d} \cdot \dfrac{f_{i,o,d,t}}{200{.}000}}_{\text{custo log\'{i}stico}}
    \end{aligned}
  }
\]

\begin{center}
  \small\color{darkgray}
  Sujeito \`{a}s restri\c{c}\~{o}es de capabilidade, balan\c{c}o de estoque, atendimento \`{a} demanda,\\
  capacidade produtiva, limite de armazenagem e n\~{a}o-negatividade.
\end{center}

% ============================================================
\section{Conjuntos (\'Indices)}

\begin{center}
\renewcommand{\arraystretch}{1.4}
\begin{tabular}{>{\bfseries\color{ballblue}}c p{8cm} c r}
  \toprule
  S\'{i}mbolo & Descri\c{c}\~{a}o & \'{I}ndice & Cardinalidade \\
  \midrule
  $I$       & SKUs (produtos finais) & $i$ & 30 \\
  $P$       & Plantas (f\'{a}bricas) & $p,\,o,\,d$ & 10 \\
  $L_p$     & Linhas da planta $p$ & $l$ & 3 por planta \\
  $L$       & Todas as linhas ($L = \bigcup_{p \in P} L_p$) & $l$ & 30 \\
  $T$       & Per\'{i}odos $(t_1, t_2, \ldots, t_{15})$ & $t$ & 15 dias \\
  $F$       & Formatos de lata do portf\'{o}lio & $-$ & 3 \\
  $\mathcal{C}$ & Pares cap\'{a}veis $(i,p,l)$: formato de $i$ = formato de $l \in L_p$ & $(i,p,l)$ & $\subseteq I \times L$ \\
  \bottomrule
\end{tabular}
\end{center}

\vspace{6pt}
\begin{mdframed}[style=destaque]
  \textbf{\color{ballblue}Conjunto $\mathcal{C}$ --- Capabilidade.}
  Os tr\^{e}s formatos s\~{a}o: \textit{9.1 Oz Sleek}, \textit{12 Oz Standard} e \textit{16 Oz Tall}.
  Cada planta $p$ (e todas as linhas $l \in L_p$) opera com exatamente 1 formato.
  Uma tripla $(i,p,l)$ pertence a $\mathcal{C}$ se o formato do SKU $i$ coincide com o
  formato da planta $p$. Fora de $\mathcal{C}$, a produ\c{c}\~{a}o \'{e} proibida ($x_{i,p,l,t} = 0$).
  Os $\approx 10\%$ de pedidos incompat\'{i}veis s\~{a}o resolvidos via transfer\^{e}ncia $f_{i,o,d,t}$
  de uma planta apta.
\end{mdframed}

% ============================================================
\section{Par\^{a}metros}

\subsection*{Par\^{a}metros log\'{i}sticos \textnormal{\small(fonte: \texttt{input\_logistic\_costs\_artificial.csv})}}

Seja o par $(o, d)$ com $o, d \in P$:

\begin{align*}
  cb_{o,d} &: \text{custo de combust\'{i}vel da rota } o \to d \\
  cp_{o,d} &: \text{custo de ped\'{a}gio da rota } o \to d \\
  cl_{o,d} &: \text{custo de m\~{a}o de obra da rota } o \to d \\
  ct_{o,d} &= cb_{o,d} + cp_{o,d} + cl_{o,d} \qquad \text{(custo total do frete por caminh\~{a}o)}
\end{align*}

\textbf{Fluxo interno:} quando $o = d$, fixamos $ct_{o,d} = 0$.

\subsection*{Par\^{a}metros de produ\c{c}\~{a}o \textnormal{\small(fontes: \texttt{input\_production\_rate} e \texttt{input\_sku\_costs})}}

\begin{align*}
  r_{p,l,t}   &: \text{taxa de produ\c{c}\~{a}o da linha } l \in L_p \text{ no dia } t \text{ (latas/hora)} \\
               &\quad \text{\small N\~{a}o indexado por SKU --- a taxa depende do formato da planta.} \\
               &\quad \text{\small Aus\^{e}ncia de registro implica produ\c{c}\~{a}o proibida naquele dia.} \\[6pt]
  k_{i,p,l}   &: \text{custo de produ\c{c}\~{a}o do SKU } i \text{ na linha } l \in L_p \text{ (R\$/lata)} \\
  v_{i,p,l}   &: \text{custo de estocagem do SKU } i \text{ na linha } l \in L_p \text{ (R\$/lata/dia)}
\end{align*}

\subsection*{Par\^{a}metros de demanda e armazenagem \textnormal{\small(fontes: \texttt{input\_production\_need} e \texttt{input\_line\_storage\_limit})}}

\begin{align*}
  \eta_{i,p,l,t}  &: \text{demanda do SKU } i \text{ alocada \`{a} linha } l \in L_p \text{ com prazo no dia } t \text{ (latas)} \\
                   &\quad \text{\small Pode ser incompat\'{i}vel com o formato da planta (ru\'{i}do de $\approx$10\%).}\\[6pt]
  e^{\max}_{p,l}  &: \text{capacidade m\'{a}xima de armazenagem da linha } l \in L_p \text{ (latas)} \\
                   &\quad \text{\small Limite f\'{i}sico fixo --- n\~{a}o varia por SKU nem por per\'{i}odo.}
\end{align*}

% ============================================================
\section{Fun\c{c}\~{a}o Auxiliar}

\[
  \phi : L \to P, \quad \phi(l) = p \qquad
  \text{(mapeamento hier\'{a}rquico: linha } l \mapsto \text{planta } p = \phi(l)\text{)}
\]

Como $L = \bigcup_{p \in P} L_p$ e $|L_p| = 3$ para todo $p$, a fun\c{c}\~{a}o $\phi$ \'{e} sobrejetora
e determina univocamente a planta de origem de qualquer produ\c{c}\~{a}o ou transfer\^{e}ncia.
Nos \'{i}ndices, sempre que aparecer $p$ e $l$ juntos, vale $p = \phi(l)$.

% ============================================================
\section{Vari\'{a}veis de Decis\~{a}o}

São cinco variáveis de decisão, todas contínuas e não-negativas, projetadas para mapear perfeitamente a física do chão de fábrica:

\vspace{8pt}
\begin{center}
\renewcommand{\arraystretch}{1.6}
\begin{tabular}{>{\bfseries\color{ballblue}}c p{9cm} l}
\toprule
Variável & Descrição & Domínio \\
\midrule
$x_{i,p,l,t}$ & Volume produzido do SKU $i$ na linha $l \in L_p$ da planta $p$, no dia $t$ (latas). & $\mathbb{R}_+$ \\[4pt]
$e_{i,p,l,t}$ & Nível de estoque do SKU $i$ na linha $l \in L_p$ ao final do dia $t$ (latas). & $\mathbb{R}_+$ \\[4pt]
$a_{i,p,l,t}$ & Volume de atendimento local do SKU $i$ a partir da linha $l \in L_p$ no dia $t$ (latas). & $\mathbb{R}_+$ \\[4pt]
$f_{i,o,d,t}$ & Volume do SKU $i$ transferido da planta $o$ para a planta $d$ no dia $t$ (latas). & $\mathbb{R}_+$ \\[4pt]
$g_{i,p,l,t}$ & Volume do SKU $i$ movimentado da doca de entrada da planta $p$ para o armazém da linha $l \in L_p$ no dia $t$ (latas). & $\mathbb{R}_+$ \\[4pt]
$s_{i,p,l,t}$ & Volume do SKU $i$ retirado do armazém da linha $l \in L_p$ e movimentado para a doca de saída da planta $p$ no dia $t$ (latas). & $\mathbb{R}_+$ \\
\bottomrule
\end{tabular}
\end{center}

\vspace{10pt}


% ── O nó de transbordo (Variáveis g e s) ──────────────────────────────────────
\begin{mdframed}[style=destaque]
\textbf{\color{ballblue}O que são as variáveis de transbordo $g_{i,p,l,t}$ e $s_{i,p,l,t}$?}

\medskip Fisicamente, os caminhões de frete ($f$) chegam e saem pelo pátio principal da fábrica (a doca da planta $p$), e não diretamente na porta de cada linha de produção $l$. 

\begin{itemize}[leftmargin=1.5em, itemsep=2pt]
\item $g_{i,p,l,t}$ \textbf{(Entrada):} representa as ``empilhadeiras''. Elas pegam as latas descarregadas pelos caminhões no pátio e as distribuem para os armazéns internos das linhas.
\item $s_{i,p,l,t}$ \textbf{(Saída):} representa as empilhadeiras fazendo o caminho inverso. Elas retiram as latas dos armazéns das linhas e as levam até a doca para encher os caminhões que vão para a estrada.
\end{itemize}

\medskip \textbf{Por que não usar $f$ direto no estoque?} 
Se não separássemos esse transbordo, a matemática exigiria que a planta de origem ditasse a organização interna do armazém da planta de destino. Pior ainda: ao somarmos o frete que chega na planta $p$ diretamente na equação do armazém da linha $l$, o solver ``clonaria'' as latas recebidas para todos os armazéns simultaneamente. As variáveis $g$ e $s$ garantem o Princípio da Separação Logística, criando um nó de transbordo perfeito.
\end{mdframed}

% ── Justificativa de f ───────────────────────────────────────────────────────
\begin{mdframed}[style=destaque]
  \textbf{\color{ballblue}Por que $f_{i,o,d,t}$ e n\~{a}o simplesmente $x_{i,p,l,t}$ com destino?}

  \medskip
  O custo de frete \'{e} pago pelo \textbf{movimento f\'{i}sico de latas entre plantas} ---
  n\~{a}o pela produ\c{c}\~{a}o em si. Um caminh\~{a}o s\'{o} \'{e} contratado quando latas
  \emph{saem} de uma planta e \emph{chegam} em outra. Se a planta de Jacare\'{i} produz
  1.000.000 de latas e guarda tudo no estoque local por 3 dias sem transferir nada,
  nenhum caminh\~{a}o \'{e} usado --- custo log\'{i}stico zero.

  \medskip
  O caminho f\'{i}sico das latas \'{e}:
  \[
    x_{i,p,l,t}
    \;\xrightarrow{\text{vai ao estoque}}\;
    e_{i,p,l,t}
    \;\xrightarrow{\text{sai da planta } p}\;
    f_{i,\,p,\,d,\,t}
    \;\xrightarrow{\text{chega em } d}\;
    \text{atendimento}
  \]

  Cada vari\'{a}vel carrega seu pr\'{o}prio custo:
  $x$ paga $k_{i,p,l}$ (produ\c{c}\~{a}o),\; $e$ paga $v_{i,p,l}$ (estocagem),\;
  $f$ paga $ct_{o,d}$ (frete). Unir tudo em uma s\'{o} vari\'{a}vel impediria o modelo
  de decidir entre ``produzir agora e transferir depois'' ou ``estocar e evitar frete''.
\end{mdframed}

\vspace{10pt}

% ── Balanço de estoque ───────────────────────────────────────────────────────
\begin{mdframed}[style=destaque]
  \textbf{\color{ballblue}Como o estoque $e_{i,p,l,t}$ se forma --- intui\c{c}\~{a}o do balan\c{c}o}

  \medskip
  O n\'{i}vel de estoque ao final do dia $t$ resulta de cinco fluxos:

  \[
    \begin{aligned}
      e_{i,p,l,t} \;=\;\;
      & \underbrace{e_{i,p,l,\,t-1}}_{\text{estoque anterior}}
        \;+\; \underbrace{x_{i,p,l,t}}_{\text{produ\c{c}\~{a}o do dia}}
        \;+\; \underbrace{\sum_{o \in P} f_{i,\,o,\,p,\,t}}_{\text{recebido de outras plantas}} \\[10pt]
      & -\; \underbrace{\sum_{d \in P} f_{i,\,p,\,d,\,t}}_{\text{enviado a outras plantas}}
        \;-\; \underbrace{a_{i,p,l,t}}_{\text{atendimento local}}
    \end{aligned}
  \]

  \smallskip
  \textit{Esta equa\c{c}\~{a}o ser\'{a} formalizada como restri\c{c}\~{a}o de igualdade no modelo.
  Aqui serve apenas para interpretar o que cada vari\'{a}vel representa no balan\c{c}o f\'{i}sico
  do armazém.}
\end{mdframed}

\vspace{10pt}

% ── Atendimento local ────────────────────────────────────────────────────────
\begin{mdframed}[style=destaque]
  \textbf{\color{ballblue}O que \'{e} o atendimento local $a_{i,p,l,t}$?}

  \medskip
  $a_{i,p,l,t}$ representa as latas do SKU $i$ que \textbf{saem do estoque da linha
  $l \in L_p$ para satisfazer a demanda $\eta_{i,p,l,t}$} no dia $t$ --- sem sair da planta.
  \'{E} uma sa\'{i}da interna: as latas foram produzidas (localmente ou recebidas via
  transfer\^{e}ncia), ficaram no armazém, e s\~{a}o consumidas para fechar o pedido daquela linha.

  \medskip
  A distin\c{c}\~{a}o entre $a_{i,p,l,t}$ e $f_{i,o,d,t}$ \'{e} fundamental:
  \begin{itemize}[leftmargin=1.5em, itemsep=2pt]
    \item $f_{i,o,d,t}$: latas que \textbf{viajam entre plantas} --- geram custo de frete.
    \item $a_{i,p,l,t}$: latas que \textbf{saem do estoque para o cliente local} --- sem frete.
  \end{itemize}

  \medskip
  O atendimento local tamb\'{e}m resolve os \textbf{10\% de pedidos com ru\'{i}do de formato}:
  a linha incompat\'{i}vel recebe $f$ de uma planta apta, armazena no estoque, e usa
  $a_{i,p,l,t}$ para atender o pedido mal alocado.
\end{mdframed}

% ============================================================
\section{Fun\c{c}\~{a}o Objetivo}

Minimizar o custo total de opera\c{c}\~{a}o da malha, somando os tr\^{e}s termos:

\[
  \begin{aligned}
    \min \; Z \;=\;\;
    & \sum_{p \in P}\;\sum_{\substack{l \in L_p \\ (i,p,l)\,\in\,\mathcal{C}}}\;\sum_{i \in I}\;\sum_{t \in T}
        k_{i,p,l} \cdot x_{i,p,l,t}
      \;+\;
      \sum_{p \in P}\sum_{l \in L_p}\sum_{i \in I}\sum_{t \in T}
        v_{i,p,l} \cdot e_{i,p,l,t} \\[10pt]
    & +\;
      \sum_{i \in I}\sum_{o \in P}\sum_{\substack{d \in P \\ d \,\neq\, o}}\sum_{t \in T}
        ct_{o,d} \cdot \frac{f_{i,o,d,t}}{200{.}000}
  \end{aligned}
\]

\vspace{8pt}
\begin{center}
\renewcommand{\arraystretch}{1.3}
\small
\begin{tabular}{>{\bfseries\color{ballblue}}c l l}
  \toprule
  Termo & F\'{o}rmula & Interpreta\c{c}\~{a}o \\
  \midrule
  Produ\c{c}\~{a}o
    & $k_{i,p,l} \cdot x_{i,p,l,t}$
    & Custo unit\'{a}rio (R\$/lata) $\times$ volume produzido \\
  Estoque
    & $v_{i,p,l} \cdot e_{i,p,l,t}$
    & Custo de carregar 1 lata no armazém por 1 dia \\
  Log\'{i}stica
    & $ct_{o,d} \cdot \dfrac{f_{i,o,d,t}}{200{.}000}$
    & Custo por caminh\~{a}o $\times$ n\'{u}mero de caminh\~{o}es \\
  \bottomrule
\end{tabular}
\end{center}

\vspace{6pt}
\begin{center}
  \small\color{darkgray}
  Capacidade do caminh\~{a}o: $25\text{ pallets} \times 8{.}000\text{ latas} = 200{.}000\text{ latas/caminh\~{a}o}$
\end{center}

% ============================================================
\newpage
\section{Restri\c{c}\~{o}es}

% ── R1 ───────────────────────────────────────────────────────────────────────
\begin{mdframed}[style=destaque]
\textbf{\color{ballblue}R1 --- Limite de armazenagem.}

\medskip
O armazém de cada linha tem uma capacidade f\'{i}sica $e^{\max}_{p,l}$ que n\~{a}o pode ser
ultrapassada em nenhum dia, independentemente do SKU. A restri\c{c}\~{a}o soma o estoque
de todos os produtos porque o espa\c{c}o \'{e} compartilhado:

\[
  \sum_{i \in I} e_{i,p,l,t} \;\leq\; e^{\max}_{p,l}
  \qquad \forall\, p \in P,\; l \in L_p,\; t \in T
\]
\end{mdframed}

\vspace{8pt}

% ── R2 ───────────────────────────────────────────────────────────────────────
\begin{mdframed}[style=destaque]
\textbf{\color{ballblue}R2 --- Balan\c{c}o de estoque.}

\medskip O n\'ivel de estoque ao final do dia $t$ \'e determinado pelo que havia ontem, pelo que a m\'aquina produziu hoje, pelo que as empilhadeiras trouxeram do p\'atio para a linha ($g$), subtraindo o que as empilhadeiras levaram para o p\'atio ($s$) e o que foi entregue diretamente ao cliente local ($a$):

$$e_{i,p,l,t} = e_{i,p,l,t-1} + x_{i,p,l,t} + g_{i,p,l,t} - s_{i,p,l,t} - a_{i,p,l,t}$$
$$\qquad \forall\, i \in I,\; p \in P,\; l \in L_p,\; t \in T$$

\medskip Esta \'e uma restri\c{c}\~ao de \textbf{igualdade} --- o balan\c{c}o f\'isico deve fechar exatamente. Ela acopla os 15 per\'iodos entre si: o estoque de hoje \'e condi\c{c}\~ao inicial de amanh\~a.
\end{mdframed}


\vspace{8pt}

% ── R3 ───────────────────────────────────────────────────────────────────────
\begin{mdframed}[style=destaque]
\textbf{\color{ballblue}R3 --- Atendimento \`{a} demanda.}

\medskip
O atendimento local $a_{i,p,l,t}$ deve cobrir integralmente a demanda $\eta_{i,p,l,t}$
no dia do prazo:

\[
  a_{i,p,l,t} \;\geq\; \eta_{i,p,l,t}
  \qquad \forall\, i \in I,\; p \in P,\; l \in L_p,\; t \in T
\]

A produ\c{c}\~{a}o e as transfer\^{e}ncias recebidas \textbf{n\~{a}o aparecem aqui explicitamente}
porque j\'{a} est\~{a}o capturadas em R2: o solver s\'{o} consegue ter $a_{i,p,l,t}$ grande
o suficiente se o balan\c{c}o de estoque (R2) tiver sido alimentado por produ\c{c}\~{a}o
local ou transfer\^{e}ncia recebida. Incluir $x$ e $f$ nesta restri\c{c}\~{a}o seria
cont\'{a}-los duas vezes.
\end{mdframed}

\vspace{8pt}

% ── R4 ───────────────────────────────────────────────────────────────────────
\begin{mdframed}[style=destaque]
\textbf{\color{ballblue}R4 --- Capacidade produtiva (turno de 24\,h).}

\medskip
A soma das horas consumidas por todos os SKUs numa linha n\~{a}o pode ultrapassar o
turno di\'{a}rio. O tempo gasto para produzir $x_{i,p,l,t}$ latas \'{e}
$x_{i,p,l,t} / r_{p,l,t}$ horas:

\[
  \sum_{i \in I} \frac{x_{i,p,l,t}}{r_{p,l,t}} \;\leq\; 24
  \qquad \forall\, p \in P,\; l \in L_p,\; t \in T
\]

Se $(i,p,l) \notin \mathcal{C}$, ent\~{a}o $x_{i,p,l,t} = 0$ e o termo correspondente
n\~{a}o contribui para a soma.
\end{mdframed}

\vspace{8pt}

% ── R5 ───────────────────────────────────────────────────────────────────────
\begin{mdframed}[style=destaque]
\textbf{\color{ballblue}R5 --- Capacidade log\'{i}stica de expedi\c{c}\~{a}o (80 caminh\~{o}es/dia).}

\medskip
Cada planta pode expedir no m\'{a}ximo 80 caminh\~{o}es por dia, somando todos os SKUs
e todos os destinos poss\'{i}veis. O limite \'{e} de \textbf{sa\'{i}da} --- o que chega
na planta n\~{a}o entra neste c\'{o}mputo:

\[
  \sum_{i \in I} \sum_{\substack{d \in P \\ d \neq p}}
  \frac{f_{i,p,d,t}}{200{.}000} \;\leq\; 80
  \qquad \forall\, p \in P,\; t \in T
\]
\end{mdframed}

\vspace{8pt}

% ── R6 ───────────────────────────────────────────────────────────────────────


\begin{mdframed}[style=destaque]
\textbf{\color{ballblue}R6 e R7 --- Conserva\c{c}\~ao de Massa no Transbordo (Docas).}

\medskip O porteiro da f\'abrica n\~ao permite a cria\c{c}\~ao nem o sumi\c{c}o de latas no p\'atio. Tudo o que os caminh\~oes descarregam deve ser exatamente igual ao que as empilhadeiras guardam (R7). E tudo o que as empilhadeiras tiram dos galp\~oes deve ser exatamente igual ao que os caminh\~oes levam embora (R8).

\textbf{R6 (Doca de Entrada):}
$$\sum_{l \in L_p} g_{i,p,l,t} = \sum_{\substack{o \in P \\ o \neq p}} f_{i,o,p,t} \qquad \forall\, i \in I,\; p \in P,\; t \in T$$

\textbf{R7 (Doca de Sa\'ida):}
$$\sum_{l \in L_p} s_{i,p,l,t} = \sum_{\substack{d \in P \\ d \neq p}} f_{i,p,d,t} \qquad \forall\, i \in I,\; p \in P,\; t \in T$$
\end{mdframed}

\begin{mdframed}[style=destaque]
\textbf{\color{ballblue}R8 --- N\~{a}o-negatividade.}

\medskip
Todas as vari\'{a}veis de decis\~{a}o s\~{a}o cont\'{i}nuas e n\~{a}o-negativas:

\[
  x_{i,p,l,t},\; e_{i,p,l,t},\; f_{i,o,d,t},\; a_{i,p,l,t} \;\geq\; 0
  \qquad \forall\, i,\, p,\, l,\, o,\, d,\, t
\]
\end{mdframed}

% ============================================================
\newpage
\section{Formula\c{c}\~{a}o Compacta do PL}

\vspace{6pt}

\begin{mdframed}[backgroundcolor=blue!5, linecolor=ballblue, linewidth=2pt, roundcorner=5pt, innertopmargin=15pt, innerbottommargin=15pt]
\textbf{min}
\vspace{-10pt}
\begin{align*}
Z \;&=\; \sum_{p \in P}\sum_{\substack{l \in L_p \\ (i,p,l)\in\mathcal{C}}}\sum_{i \in I}\sum_{t \in T} k_{i,p,l} \cdot x_{i,p,l,t} 
\;+\; \sum_{p \in P}\sum_{l \in L_p}\sum_{i \in I}\sum_{t \in T} v_{i,p,l} \cdot e_{i,p,l,t} \\[10pt]
&\quad +\; \sum_{i \in I}\sum_{o \in P}\sum_{\substack{d \in P \\ d \neq o}}\sum_{t \in T} ct_{o,d} \cdot \frac{f_{i,o,d,t}}{200{.}000}
\end{align*}

\textbf{s.a.}
\vspace{-5pt}
\begin{align}
& \sum_{i \in I} e_{i,p,l,t} \leq e^{\max}_{p,l} && \forall\, p \in P,\; l \in L_p,\; t \in T \tag{R1} \\[10pt]
& e_{i,p,l,t} = e_{i,p,l,t-1} + x_{i,p,l,t} + g_{i,p,l,t} - s_{i,p,l,t} - a_{i,p,l,t} && \forall\, i \in I,\; p \in P,\; l \in L_p,\; t \in T \tag{R2} \\[10pt]
& a_{i,p,l,t} \geq \eta_{i,p,l,t} && \forall\, i \in I,\; p \in P,\; l \in L_p,\; t \in T \tag{R3} \\[10pt]
& \sum_{i \in I} \frac{x_{i,p,l,t}}{r_{p,l,t}} \leq 24 && \forall\, p \in P,\; l \in L_p,\; t \in T \tag{R4} \\[10pt]
& \sum_{i \in I} \sum_{\substack{d \in P \\ d \neq p}} \frac{f_{i,p,d,t}}{200{.}000} \leq 80 && \forall\, p \in P,\; t \in T \tag{R5} \\[10pt]
& \sum_{l \in L_p} g_{i,p,l,t} = \sum_{\substack{o \in P \\ o \neq p}} f_{i,o,p,t} && \forall\, i \in I,\; p \in P,\; t \in T \tag{R6} \\[10pt]
& \sum_{l \in L_p} s_{i,p,l,t} = \sum_{\substack{d \in P \\ d \neq p}} f_{i,p,d,t} && \forall\, i \in I,\; p \in P,\; t \in T \tag{R7} \\[10pt]
& x_{i,p,l,t},\; e_{i,p,l,t},\; a_{i,p,l,t},\; f_{i,o,d,t},\; g_{i,p,l,t},\; s_{i,p,l,t} \geq 0 && \forall\, i, p, l, o, d, t \tag{R8}
\end{align}
\end{mdframed}

\vspace{10pt}
{\color{accentgray}\rule{\linewidth}{0.5pt}}
\begin{center}
  \small\color{darkgray}
  Dimens\~{a}o do problema: $|I| \times |P| \times |L_p| \times |T|
  = 30 \times 10 \times 3 \times 15 = \mathbf{13{.}500}$ vari\'{a}veis base\,
  $+\; |I| \times |P|^2 \times |T| = 30 \times 100 \times 15 = \mathbf{45{.}000}$ vari\'{a}veis de transfer\^{e}ncia.
\end{center}

\end{document}